\input{_preamble_ftn.tex}
\begin{document}
\hd{FTN P1 — tematski listovi (zadaci)}{10 slotova · nivo hard/rare · rešenja u zasebnom fajlu}

\section*{Slot 1. Kompleksni brojevi}
\begin{enumerate}[label=\arabic*.]
\item Dat je kompleksan broj. Odrediti njegov moduo, glavni argument $\arg z\in(-\pi,\pi]$ i napisati broj u trigonometrijskom obliku:
\begin{enumerate}[label=\alph*)]
\item $z=-3+3\sqrt{3}\,i$;
\item $z=\dfrac{1+i}{\sqrt{3}-i}$.
\end{enumerate}

\item Koristeći Moavrovu formulu, izračunati:
\begin{enumerate}[label=\alph*)]
\item $\left(\dfrac{\sqrt{3}+i}{2}\right)^{2026}$;
\item $(1-i)^{2026}$;
\item naći sve korene četvrtog stepena broja $w=-8+8\sqrt{3}\,i$ i napisati ih u algebarskom obliku.
\end{enumerate}

\item Neka je $z=3-2i$.
\begin{enumerate}[label=\alph*)]
\item Izračunati $z\cdot\bar z$, $\dfrac{\bar z}{z}$ i $\dfrac{1}{z}$.
\item Naći realni i imaginarni deo broja $w=\dfrac{4+7i}{2+i}+\dfrac{1-3i}{1+i}.$
\end{enumerate}

\item \begin{enumerate}[label=\alph*)]
\item Rešiti jednačinu $z^2-(3+i)z+(2+2i)=0$ nad skupom kompleksnih brojeva.
\item Odrediti realne brojeve $a$ i $b$ za koje je broj $z=1+2i$ koren jednačine $z^2+(a-i)z+(b+3i)=0.$
\end{enumerate}

\item Za svako realno $t$ dat je broj $z=(t+2)+(t-4)i.$ Naći vrednost $t$ za koju je moduo $|z|$ najmanji, i izračunati odgovarajuću vrednost $|z|$ i sam broj $z$.
\end{enumerate}
\clearpage
\section*{Slot 2. Kvadratna jednačina i Vijetove formule}
\begin{enumerate}[label=\arabic*.]
\item Data je jednačina $x^{2}-mx+3=0$, $m\in\mathbb{R}$, sa korenima $x_1,x_2$. Po Vijetovim formulama $x_1+x_2=m$, $x_1x_2=3$.
\begin{enumerate}[label=\alph*)]
\item Naći sve $m$ za koje je $\dfrac{1}{x_1}+\dfrac{1}{x_2}=\dfrac{5}{3}$.
\item Naći sve $m$ za koje je $x_1^{3}+x_2^{3}=28$.
\item Naći sve $m$ za koje je $\dfrac{1}{x_1^{3}}+\dfrac{1}{x_2^{3}}=6$.
\end{enumerate}
U svim tačkama proveriti da su koreni realni.

\item \begin{enumerate}[label=\alph*)]
\item Data je jednačina $x^{2}-(m-2)x+(m-3)=0$, $m\in\mathbb{R}$. Za koje $m$ jednačina ima dvostruki (jednak) koren? Naći taj koren.
\item Za istu jednačinu naći sve $m$ za koje je jedan od korena jednak nuli, i navesti drugi koren.
\item Data je jednačina $x^{2}-2(m-2)x+(m^{2}-4)=0$, $m\in\mathbb{R}$. Odrediti sve $m$ za koje su koreni realni i: (i) imaju različite znake; (ii) oba su negativna.
\end{enumerate}

\item \begin{enumerate}[label=\alph*)]
\item Data je jednačina $x^{2}-mx+12=0$. Naći sve $m$ za koje je jedan koren tri puta veći od drugog ($x_1=3x_2$).
\item Data je jednačina $x^{2}-mx+8=0$. Naći $m$ za koje je jedan koren jednak kvadratu drugog ($x_1=x_2^{2}$).
\item Data je jednačina $x^{2}-5x+m=0$ sa korenima $x_1,x_2$. Naći $m$ za koje je $2x_1-x_2=7$.
\end{enumerate}

\item \begin{enumerate}[label=\alph*)]
\item Data je jednačina $x^{2}-mx+(m^{2}-3m)=0$, $m\in\mathbb{R}$. Na skupu onih $m$ za koje su koreni realni, naći najveću vrednost izraza $x_1^{2}+x_2^{2}$ i odgovarajuće $m$.
\item Data je jednačina $x^{2}-2mx+(3m+4)=0$, $m\in\mathbb{R}$. Na skupu onih $m$ za koje su koreni realni, naći najmanju vrednost izraza $x_1^{2}+x_2^{2}$ i odgovarajuće $m$.
\end{enumerate}

\item \begin{enumerate}[label=\alph*)]
\item Tri broja čine aritmetičku progresiju. Njihov zbir je $15$, a zbir njihovih kvadrata je $83$. Naći te brojeve.
\item Tri pozitivna broja čine geometrijsku progresiju. Njihov zbir je $26$, a proizvod je $216$. Naći te brojeve.
\item Rešiti jednačinu $|x-1|+|2x-4|=5$.
\end{enumerate}
\end{enumerate}
\clearpage
\section*{Slot 3. Logaritmi}
\begin{enumerate}[label=\arabic*.]
\item
\begin{enumerate}[label=\alph*)]
\item Naći oblast definisanosti funkcije $f(x)=\log_{1/3}\left(x^{2}-2x\right)$.
\item Rešiti nejednačinu $\log_{1/3}\left(x^{2}-2x\right) > -1$.
\end{enumerate}

\item Rešiti jednačine:
\begin{enumerate}[label=\alph*)]
\item $\log^{2} x-\log\left(x^{3}\right)=4$ (dekadni logaritam);
\item $\log_{2} x+\log_{x} 2=\dfrac{5}{2}$.
\end{enumerate}

\item Neka je $a=\log_{2}3$ i $b=\log_{2}5$. Izraziti pomoću $a$ i $b$:
\begin{enumerate}[label=\alph*)]
\item $\log_{12}150$;
\item $\log_{5}0{,}24$.
\end{enumerate}

\item
\begin{enumerate}[label=\alph*)]
\item Rešiti jednačinu sa promenljivom osnovom $\log_{x-1}(2x-2)=2$.
\item Za koju vrednost $x$ su brojevi $\log_{3}2,\ \log_{3}\left(2^{x}-1\right),\ \log_{3}\left(2^{x}+3\right)$ uzastopni članovi aritmetičke progresije?
\end{enumerate}

\item Rešiti sisteme jednačina:
\begin{enumerate}[label=\alph*)]
\item $\begin{cases}\log_{2}x+\log_{2}y=5,\\[2pt] \log_{2}x-\log_{2}y=1;\end{cases}$
\item $\begin{cases}\log x+\log y=\log 20,\\[2pt] \log x-\log y=\log 5\end{cases}$ (dekadni logaritam).
\end{enumerate}
\end{enumerate}
\clearpage
\section*{Slot 4. Eksponencijalne jednačine i nejednačine}
\begin{enumerate}[label=\arabic*.]
\item Rešiti:
\begin{enumerate}[label=\alph*)]
\item $4^{x}-3\cdot 2^{x}-4=0$;
\item $9^{x}-4\cdot 3^{x}+3\le 0$.
\end{enumerate}

\item Rešiti:
\begin{enumerate}[label=\alph*)]
\item $4\cdot 9^{x}-13\cdot 6^{x}+9\cdot 4^{x}=0$;
\item $6\cdot 9^{x}-13\cdot 6^{x}+6\cdot 4^{x}=0$.
\end{enumerate}

\item Rešiti:
\begin{enumerate}[label=\alph*)]
\item $5^{\,x^{2}-3}=25^{\,|x|}$;
\item $2^{\,|x^{2}-2x|}\le 8$.
\end{enumerate}

\item Rešiti:
\begin{enumerate}[label=\alph*)]
\item $2^{\,x-1}=2^{\sqrt{x+1}}$;
\item $4^{\sqrt{x}}-5\cdot 2^{\sqrt{x}}+4=0$.
\end{enumerate}

\item Rešiti:
\begin{enumerate}[label=\alph*)]
\item $\log_{2}\big(\log_{3}(9^{x}-6)\big)=0$;
\item $\log_{5}\big(25^{x}-20\big)=x$.
\end{enumerate}
\end{enumerate}
\clearpage
\section*{Slot 5. Trigonometrija na intervalu}
\begin{enumerate}[label=\arabic*.]
\item Rešiti jednačine na naznačenim intervalima.
\begin{enumerate}[label=\alph*)]
\item $\cos 2x + 3\cos x + 2 = 0$ na intervalu $[0,\,2\pi)$;
\item $\cos 3x + \cos x = \cos 2x$ na intervalu $(0,\,\pi)$.
\end{enumerate}

\item Rešiti jednačine na intervalu $[0,\,2\pi)$.
\begin{enumerate}[label=\alph*)]
\item $\sin^3 x + \cos^3 x = 1$;
\item $\sin 2x = \sin x$.
\end{enumerate}

\item Rešiti jednačine na intervalu $[0,\,2\pi)$, prethodno navodeći oblast definisanosti.
\begin{enumerate}[label=\alph*)]
\item $\operatorname{tg} x + \operatorname{ctg} x = 2$;
\item $\operatorname{tg} x = 3\,\operatorname{ctg} x$.
\end{enumerate}

\item Naći sva rešenja jednačine
\[2\sin\!\left(x-\frac{\pi}{6}\right)=\sqrt3,\]
koja pripadaju intervalu $[0,\,2\pi)$, i posebno ona od njih koja leže u $\left(\dfrac{2\pi}{3},\,2\pi\right)$.

\item Data je funkcija $f(x)=2\cos x-\cos 2x-1$, posmatrana na intervalu $(-\pi,\,\pi]$.
\begin{enumerate}[label=\alph*)]
\item Naći sve nule funkcije $f$ na $(-\pi,\,\pi]$.
\item Rešiti nejednačinu $f(x)>0$ na $(-\pi,\,\pi]$.
\item Rešiti nejednačinu $f(x)<0$ na $(-\pi,\,\pi]$.
\end{enumerate}
\end{enumerate}
\clearpage
\section*{Slot 6. Vektori i analitička geometrija}
\begin{enumerate}[label=\arabic*.]

\item Dati su nekolinearni vektori $m$ i $n$, pri čemu je $|m|=2$, $|n|=3$, a ugao između njih je $\dfrac{\pi}{3}$. Stavimo $a=2m-n$ i $b=m+3n$.
\begin{enumerate}[label=\alph*)]
\item Naći dužinu vektora $a$.
\item Izračunati površinu paralelograma konstruisanog nad vektorima $a$ i $b$.
\item Za koju vrednost $t$ je vektor $m+tn$ ortogonalan na vektor $a=2m-n$?
\item Razložiti vektor $m$ po bazi $\{a,b\}$, to jest naći brojeve $x,y$ za koje je $m=xa+yb$.
\end{enumerate}

\item U prostoru su date tačke $A(1,1,1)$, $B(-1,2,2)$ i $C(0,0,3)$.
\begin{enumerate}[label=\alph*)]
\item Naći ugao $\angle BAC$.
\item Izračunati površinu trougla $ABC$.
\item Naći koordinate težišta (presečne tačke težišnih linija) trougla $ABC$.
\end{enumerate}

\item \begin{enumerate}[label=\alph*)]
\item Dijagonale paralelograma date su vektorima $d_1=(1,1,2)$ i $d_2=(3,-3,0)$. Naći površinu paralelograma.
\item Za koju vrednost parametra $\lambda$ su vektori $u=(1,2,-1)$, $v=(2,-1,3)$ i $w=(\lambda,1,2)$ komplanarni?
\item Za koju vrednost parametra $k$ je vektor $p=(2,-3,k)$ kolinearan vektoru $q=(-4,6,10)$?
\end{enumerate}

\item \begin{enumerate}[label=\alph*)]
\item Vektori $p$ i $q$ zadovoljavaju $|p|=1$, $|q|=\sqrt2$, a ugao između njih je $\dfrac{\pi}{4}$. Naći $|p+q|$ i vrednost $t$ za koju je vektor $p+tq$ ortogonalan na vektor $q$.
\item Za koju vrednost $\alpha$ su vektori $a=(\alpha,-1,2)$ i $b=(3,\alpha,-2)$ ortogonalni?
\item Za koju vrednost $\alpha$ vektori $u=(\alpha-1,3,0)$ i $v=(\alpha+1,2,1)$ imaju jednake dužine?
\end{enumerate}

\item \begin{enumerate}[label=\alph*)]
\item Dat je trougao sa temenima $A(1,1)$, $B(7,1)$, $C(3,5)$. Naći koordinate njegovog ortocentra (presečne tačke visina).
\item Za isti trougao naći centar i poluprečnik opisane kružnice i napisati njenu jednačinu.
\item Dokazati da je četvorougao sa temenima $P(0,0)$, $Q(3,4)$, $R(8,4)$, $S(5,0)$ romb.
\end{enumerate}

\end{enumerate}
\clearpage
\section*{Slot 7. Planimetrija}
\begin{enumerate}[label=\arabic*.]

\item Osnovice trapeza su $a=20$ i $c=8$. Ugao na većoj osnovici sa strane jednog kraka je $30^\circ$, a sa strane drugog kraka je $45^\circ$.
\begin{enumerate}[label=\alph*)]
\item Naći visinu trapeza i oba kraka.
\item Naći površinu trapeza.
\end{enumerate}

\item \begin{enumerate}[label=\alph*)]
\item U jednakokraki trapez sa osnovicama $a=18$ i $c=8$ može se upisati kružnica. Naći krak, poluprečnik upisane kružnice i površinu trapeza.
\item U jednakokrakom trapezu dijagonale su uzajamno normalne, a osnovice su $a=14$ i $c=6$. Naći visinu i površinu trapeza.
\end{enumerate}

\item \begin{enumerate}[label=\alph*)]
\item Razlika dijagonala romba je $14$, a stranica je $13$. Naći dijagonale, površinu i poluprečnik upisane kružnice.
\item Dijagonale romba odnose se kao $3:4$, a njegova površina je $96$. Naći dijagonale, stranicu i poluprečnik upisane kružnice.
\end{enumerate}

\item Stranice trougla čine aritmetičku progresiju, njegov obim je $30$, a jedan od uglova je $120^\circ$.
\begin{enumerate}[label=\alph*)]
\item Naći stranice trougla i njegovu površinu.
\item Naći poluprečnike upisane i opisane kružnice.
\end{enumerate}

\item Dve kružnice sa poluprečnicima $R_1=9$ i $R_2=4$ dodiruju se spolja. Na njih je povučena zajednička spoljašnja tangenta.
\begin{enumerate}[label=\alph*)]
\item Naći dužinu odsečka zajedničke spoljašnje tangente (između tačaka dodira) i površinu trapeza obrazovanog poluprečnicima u tačkama dodira, odsečkom tangente i linijom centara.
\item Naći poluprečnik kružnice koja dodiruje obe date kružnice i njihovu zajedničku tangentu (upisana u krivolinijski „ugao” između njih).
\end{enumerate}

\item U kružnici poluprečnika $R=6$ povučena je tetiva, udaljena od centra na rastojanje $3$.
\begin{enumerate}[label=\alph*)]
\item Naći dužinu tetive, centralni ugao nad kojim se ona nalazi, i periferijski ugao nad istom tetivom sa strane većeg luka.
\item Ta tetiva deli krug na dva odsečka. Koji deo (u procentima) površine kruga čini manji odsečak?
\end{enumerate}

\end{enumerate}
\clearpage
\section*{Slot 8. Stereometrija}
\begin{enumerate}[label=\arabic*.]
\item Data je prava kružna kupa, kod koje poluprečnik osnove $r$, visina $H$ i izvodnica $s$ čine (upravo tim redom) aritmetičku progresiju, pri čemu je $r=3$.
\begin{enumerate}[label=\alph*)]
\item Naći $H$ i $s$, a zatim zapreminu i površinu kupe.
\item U kupu je upisana lopta (dodiruje osnovu i omotač). Naći poluprečnik $\rho$ te lopte.
\item Izračunati odnos zapremine upisane lopte prema zapremini kupe.
\end{enumerate}

\item U pravilnu četvorostranu piramidu sa stranicom osnove $a$ i visinom $H$ upisana je kocka: jedna stranica kocke leži u osnovi piramide (centrirana), a četiri gornja temena kocke leže na bočnim stranama piramide.
\begin{enumerate}[label=\alph*)]
\item Izraziti ivicu kocke $b$ kao funkciju od $a$ i $H$.
\item Za $a=6$, $H=12$ naći ivicu kocke, njenu zapreminu i površinu.
\item Za iste $a,H$ naći odnos zapremine kocke prema zapremini piramide.
\end{enumerate}

\item Pravougli trougao sa katetama $3$ i $4$ smešten je u ravan tako da je kateta dužine $4$ paralelna pravoj $\ell$, a sam trougao u celini leži sa jedne strane prave $\ell$; najbliža pravoj $\ell$ kateta (dužine $4$) udaljena je od $\ell$ na rastojanje $2$. Trougao rotira oko prave $\ell$.
\begin{enumerate}[label=\alph*)]
\item Opisati telo koje se dobija i naći njegovu zapreminu.
\item Naći površinu tela nastalog rotacijom.
\end{enumerate}
\emph{Uputstvo: mogu se koristiti Guldinove (Papusove) teoreme ili metoda preseka.}

\item U pravu kružnu kupu sa poluprečnikom osnove $R=6$ i visinom $H=9$ upisan je valjak: donja osnova valjka leži u osnovi kupe, gornja kružnica — na omotaču kupe. Neka je $x$ — poluprečnik valjka.
\begin{enumerate}[label=\alph*)]
\item Izraziti visinu valjka i njegovu zapreminu $V(x)$ kao funkciju od $x$ (navesti oblast definisanosti).
\item Pomoću izvoda naći $x$ za koje je zapremina maksimalna, i izračunati tu maksimalnu zapreminu.
\item Naći odnos maksimalne zapremine valjka prema zapremini kupe.
\end{enumerate}

\item U zarubljenu kupu sa poluprečnicima osnova $R=4$ i $r=1$ upisana je lopta (dodiruje obe osnove i omotač).
\begin{enumerate}[label=\alph*)]
\item Dokazati da je visina zarubljene kupe jednaka $h=2\sqrt{Rr}$, i naći $h$, izvodnicu $s$ i poluprečnik $\rho$ lopte.
\item Naći zapremine lopte i zarubljene kupe i njihov odnos.
\item Naći površinu zarubljene kupe, površinu lopte i njihov odnos. Uporediti sa odnosom zapremina.
\end{enumerate}

\item U pravilnoj četvorostranoj piramidi sa osnovom $ABCD$ (kvadrat) i vrhom $S$ dve naspramne bočne strane su uzajamno normalne (diedar između strana $SAB$ i $SCD$ je $90^\circ$). Stranica osnove je $a=6$.
\begin{enumerate}[label=\alph*)]
\item Dokazati da je visina piramide jednaka $H=\dfrac{a}{2}$, i naći zapreminu.
\item Naći apotemu (visinu bočne strane) i površinu piramide.
\item Naći diedar pri ivici osnove (između bočne strane i osnove).
\end{enumerate}

\item Data je kocka kod koje je površina (brojno) jednaka zapremini.
\begin{enumerate}[label=\alph*)]
\item Naći ivicu kocke $a$.
\item Naći rastojanje od temena kocke do ravni koja prolazi kroz tri njemu susedna temena (ravan koja odseca ugao kocke).
\item Naći odnos zapremine lopte upisane u kocku prema zapremini lopte opisane oko kocke.
\end{enumerate}
\end{enumerate}
\clearpage
\section*{Slot 9. Ispitivanje funkcije i integral}
\begin{enumerate}[label=\arabic*.]
  \item Data je funkcija $f(x)=\dfrac{x^2-3}{x-2}$.
  \begin{enumerate}[label=\alph*)]
    \item Naći oblast definisanosti, vertikalnu i kosu asimptotu.
    \item Naći intervale monotonosti i ekstremne vrednosti.
    \item Izračunati $\displaystyle\int_{3}^{4} f(x)\,dx$.
  \end{enumerate}

  \item Data je funkcija $f(x)=x^2 e^{-x}$.
  \begin{enumerate}[label=\alph*)]
    \item Naći intervale monotonosti i ekstremne vrednosti.
    \item Izračunati $\displaystyle\lim_{x\to+\infty} f(x)$, $\displaystyle\lim_{x\to+\infty}\frac{f(x)}{x}$, $\displaystyle\lim_{x\to+\infty}\frac{f(x)}{x^2}$ i navesti horizontalnu asimptotu kada $x\to+\infty$.
    \item Izračunati $\displaystyle\int_{0}^{1} x^2 e^{-x}\,dx$.
  \end{enumerate}

  \item \begin{enumerate}[label=\alph*)]
    \item Za funkciju $g(x)=(x^2-4)^3$ naći sve tačke u kojima je $g'(x)=0$, i odrediti u kojima od njih postoji ekstremna vrednost, a u kojima ne (sa obrazloženjem).
    \item Izračunati smenom $\displaystyle\int_{2}^{3} x\,(x^2-4)^2\,dx$.
    \item Naći površinu oblasti ograničene grafikom funkcije $h(x)=x^3-4x$ i $Ox$ osom na odsečku $[-2,2]$.
  \end{enumerate}

  \item Data je funkcija $f(x)=\dfrac{x^2+2}{x+1}$.
  \begin{enumerate}[label=\alph*)]
    \item Naći oblast definisanosti, vertikalnu i kosu asimptotu.
    \item Napisati jednačine tangente i normale na grafik u tački sa apscisom $x_0=0$.
    \item Izračunati $\displaystyle\lim_{x\to+\infty}\frac{f(x)}{x}$.
  \end{enumerate}

  \item \begin{enumerate}[label=\alph*)]
    \item Za funkciju $f(x)=\dfrac{3x-2}{x-1}$ naći oblast definisanosti i sve asimptote (vertikalnu i horizontalnu).
    \item Izračunati $\displaystyle\lim_{x\to 0}\frac{e^{2x}-1}{x}$ i $\displaystyle\lim_{x\to+\infty} x\sin\frac{3}{x}$.
    \item Izračunati $\displaystyle\lim_{x\to+\infty}\bigl(\sqrt{x^2+x}-x\bigr)$.
  \end{enumerate}

  \item Izračunati određene integrale:
  \begin{enumerate}[label=\alph*)]
    \item $\displaystyle\int_{1}^{e} x\ln x\,dx$ \;(parcijalna integracija);
    \item $\displaystyle\int_{0}^{1} x\sqrt{1+x^2}\,dx$ \;(smenom);
    \item $\displaystyle\int_{1}^{e} \frac{\ln x}{x}\,dx$ \;(smenom) \; i \; $\displaystyle\int_{0}^{\pi/2} x\sin x\,dx$ \;(parcijalna integracija).
  \end{enumerate}
\end{enumerate}
\clearpage
\section*{Slot 10. Kombinatorika}
\begin{enumerate}[label=\arabic*.]
\item Ima $6$ kuglica koje treba rasporediti u $3$ kutije. Na koliko načina se to može učiniti u svakom od slučajeva:
\begin{enumerate}[label=\alph*)]
\item kuglice su različite, kutije su različite (bez ograničenja);
\item kuglice su različite, kutije su različite, ali svaka kutija mora biti neprazna;
\item kuglice su jednake, kutije su različite, bez ograničenja;
\item kuglice su različite, kutije su jednake, svaka kutija je neprazna; zatim naći broj načina ako su i kuglice jednake (svaka kutija neprazna).
\end{enumerate}

\item Naći broj celobrojnih rešenja jednačine
\[
x_1+x_2+x_3+x_4=20
\]
pod uslovima:
\begin{enumerate}[label=\alph*)]
\item $x_i\ge 0$ za sve $i$;
\item $x_i\ge 1$ za sve $i$;
\item $x_i\ge 2$ za sve $i$.
\end{enumerate}

\item Od cifara skupa $\{0,1,2,3,4,5\}$ obrazuju se \emph{četvorocifreni} brojevi \textbf{sa različitim} ciframa (cifre se u broju ne ponavljaju). Koliko je takvih brojeva:
\begin{enumerate}[label=\alph*)]
\item ukupno;
\item deljivih sa $5$;
\item parnih.
\end{enumerate}

\item \begin{enumerate}[label=\alph*)]
\item Za okrugli sto treba posaditi $4$ mladića i $4$ devojke tako da se mladići i devojke \textbf{naizmenično} smenjuju. Na koliko načina se to može učiniti (rasporedi koji se poklapaju pri obrtanju stola smatraju se jednakim)?
\item U red od $8$ mesta treba posaditi $8$ ljudi, među kojima $3$ određene osobe ne smeju sedeti zajedno (nikoje dve od te trojice ne zauzimaju susedna mesta). Na koliko načina se to može učiniti?
\end{enumerate}

\item Razmatraju se sve moguće permutacije slova reči \textbf{MATEMATIKA} (slova: $M,M,A,A,A,T,T,E,I,K$).
\begin{enumerate}[label=\alph*)]
\item Koliko ukupno ima različitih permutacija?
\item U koliko od njih sva tri slova $A$ stoje zajedno?
\item U koliko od njih nikoje dve slova $T$ ne stoje zajedno?
\end{enumerate}

\item \begin{enumerate}[label=\alph*)]
\item Iz grupe u kojoj je $7$ mladića i $5$ devojaka bira se komisija od $4$ osobe. Na koliko načina se to može učiniti tako da u komisiji bude \textbf{bar jedna} devojka?
\item Na koliko načina se broj $6$ može napisati u obliku zbira $3$ prirodna sabirka, \emph{ako se uzima u obzir redosled sabiraka} (kompozicije)? A na koliko — ako se redosled \emph{ne uzima u obzir} (razbijanja)?
\item Kod se sastoji od dva \textbf{različita} slova iz skupa $\{A,B,C,D,E,F\}$, za kojima slede tri cifre, pri čemu prva cifra nije $0$ (na ostale cifre nema ograničenja, cifre se mogu ponavljati). Koliko se različitih kodova može sastaviti?
\end{enumerate}
\end{enumerate}

\end{document}
